\documentclass[12pt,letterpaper]{article}
\usepackage[utf8]{inputenc}
\usepackage[spanish]{babel}
\usepackage{amsmath}
\usepackage{amsfonts}
\usepackage{amssymb}
\usepackage[margin=2.5cm]{geometry}

\begin{document}

\begin{center}
    \Large \textbf{Evaluación de Examen 2: Unidad II} \\
    \large Física para Ciencias de la Salud (005-1713) \\
    \vspace{0.5cm}
    \textbf{Estudiante:} Glanyernys Pérez \quad \textbf{C.I.:} 33.171.611
    \vspace{0.3cm} \\
    \large \textbf{Calificación Final: 8/10 puntos}
\end{center}

\vspace{0.5cm}

\section*{Análisis de Resultados}

\subsection*{1. Cinemática (Aceleración media)}
\textbf{Procedimiento y resultado:} Extrae de manera impecable los datos iniciales y plantea la resta de velocidades dividida entre el tiempo ($\frac{0,6 \text{ m/s} - 0 \text{ m/s}}{0,15 \text{ s}}$). La división es correcta y el resultado analítico es exacto: $4 \text{ m/s}^2$. \\
\textbf{Veredicto:} \textbf{Correcto} (2/2 pts).

\subsection*{2. Dinámica (Leyes de Newton)}
\textbf{Procedimiento y resultado:} Plantea correctamente el razonamiento matemático (dividiendo los $150 \text{ N}$ de fuerza entre los $25 \text{ kg}$ de masa). El cálculo aritmético es exacto ($6$); no obstante, comete un error en el análisis dimensional de la respuesta final al expresar las unidades de aceleración al cubo ($6 \text{ m/s}^3$) en lugar de al cuadrado ($\text{m/s}^2$). \\
\textbf{Veredicto:} \textbf{Parcialmente correcto} (Penalización leve por error en la unidad fundamental) (1,5/2 pts).

\subsection*{3. Trabajo Mecánico}
\textbf{Procedimiento y resultado:} Identifica y plasma la fórmula base del trabajo mecánico, multiplicando la fuerza ejercida de $400 \text{ N}$ por la distancia de $0,8 \text{ m}$. El resultado es matemáticamente correcto: $320 \text{ J}$. \\
\textbf{Veredicto:} \textbf{Correcto} (2/2 pts).

\subsection*{4. Energía Potencial}
\textbf{Procedimiento y resultado:} La estudiante extrae los datos del problema de manera estructurada (masa, altura y gravedad), pero detiene el desarrollo en este punto. Deja el resto de la página en blanco, omitiendo el planteamiento de la fórmula, la sustitución numérica y el resultado final ($17,64 \text{ J}$). \\
\textbf{Veredicto:} \textbf{Incompleto} (Se otorgan puntos parciales únicamente por la correcta extracción de los datos) (0,5/2 pts).

\subsection*{5. Potencia Mecánica}
\textbf{Procedimiento y resultado:} Relaciona correctamente el trabajo y el tiempo planteando la fracción respectiva ($\frac{75 \text{ J}}{60 \text{ s}}$). El desarrollo de la división arroja la cifra correcta y verificada de $1,25 \text{ W}$. \\
\textbf{Veredicto:} \textbf{Correcto} (2/2 pts).

\vspace{0.5cm}
\section*{Comentarios Generales y Sugerencias}
El examen evidencia un buen dominio general de las operaciones matemáticas básicas de la Unidad II y presenta una muy buena estructura de resolución.
\begin{itemize}
    \item Se recomienda tener especial \textbf{cuidado con el análisis dimensional}. La aceleración en el Sistema Internacional (S.I.) siempre lleva la unidad de tiempo al cuadrado ($\text{m/s}^2$). Además, los símbolos del S.I. distinguen mayúsculas y minúsculas: kilogramos es $\text{kg}$ (no $\text{Kg}$), y Watts es $\text{W}$ (no $\text{w}$).
    \item Vigilar la **gestión del tiempo** en las evaluaciones para evitar que un problema que se inició correctamente (como la extracción de datos en la pregunta 4) quede completamente sin resolver.
    \item ¡Buen trabajo! Con pulir estos pequeños detalles teóricos de las unidades alcanzarás la excelencia en la próxima unidad.
\end{itemize}

\end{document}